\newpage
\section{P. Svidler - M. Carlsen, Norway Chess 2014}
\begin{multicols}{2}
    \chessgameinfo{Norway Chess}{P. Svidler}{M. Carlsen}{}{2014}{1/2-1/2}
    \newchessgame
    \mainline[level=1]{
        1.c4 e5 2.Nc3 Nc6 3.Nf3 f5}
    
    \mythoughts{
        Black is playing the Grand Prix Attack. The idea is f5 and \symbishop b4xc3.
        He castles kingside, brings the queen to h5 via e8, and the knight to g4.
        That yields three attackers plus a pawn against the king. An f4 break could allow
        a rook sacrifice on f3 and mate on h2.

        Of course this is only a sketch; the same idea can crop up in other openings.
        White could adopt a Rossolimo against the Sicilian (\variation[invar]{1. e4 c5 2. Nf3 Nc6 3. Bb5})
        and mirror the whole plan with reversed colors.
    }
    \mainline[level=1]{4.d3 Nf6 5.g3 Bb4 6.Bg2 Bxc3+ 7.bxc3}
    
    \chessboard
    \mythoughts{
        Black has already traded the knight on c3 for the bishop, as planned, and still intends to
        castle kingside.

        Where should the light-squared bishop go? A fianchetto is wrong: b7 is loose, and on a6
        the bishop has too little to do.

        I do not want to let White play \symbishop a3 and spoil my castling, so I castle now.
    }

    \mainline[level=1]{7...d6}
    
    \myreflection{
        \variation[invar]{7...O-O} allows \variation[]{8. c5}, leaving a weakness on d6 for Black.
        My thoughts above are not entirely wrong, but as Lasker put it, ``if you find
        a good move, look for a better one.'' \variation[]{7... d6} is better.
    }
    \mainline[level=1]{8.O-O}

    \mythoughts{
        As planned, I castle on the kingside.
    }

    \mainline[level=1]{8...O-O 9.Rb1}
    \mythoughts{
        I bring the queen to e8 as planned.
    }

    \mainline[level=1]{9...Qe8}
    \myreflection{
        Magnus has followed the script I expected. Over the board, play is not only
        move by move but plan by plan: a plan is a short sequence of moves.
        As long as nothing forces a change of course, following the plan saves time.
    }
    \mainline[level=1]{10.Qb3}
    
    \mythoughts{
        I dislike having my queen on the same diagonal as White's queen, so I would
        play \variation[invar]{10... Kh8}.
    }
    \mainline[level=1]{10...b6}
    \myreflection{
        \variation[invar]{10... Kh8} is playable. \variation{11. c5+ Be6} White loses a pawn.
        Judging how much prophylaxis you need is hard. If the king stays on g8, Black must watch
        the c5 break for the rest of the game. Former world champion Petrosian would tuck the king on h8 without hesitation.

        Peter Leko is often on commentary. He likes to say ``I don't want adventure. I don't want
        to calculate,'' while a younger co-commentator rattles off concrete lines to show that all is well.

        Am I getting old?
    }

    \mainline[level=1]{11.Nh4}
    
    \mythoughts{
        In closed systems where White has a light-squared fianchetto, Black often struggles to
        develop on the queenside because of the pin. With White's queen on b3, I should use the chance
        to bring out the knight and rook.
    }
    \mainline[level=1]{11...Na5 12.Qa3}
    
    \mythoughts{
        We stick to the plan and develop the rook to a safer square.
    }
    \mainline[level=1]{12...Rb8 13.Be3}
    
    \chessboard

    \mythoughts{
        Black is aiming for ...f4.
        White has no time for \variation[invar]{14. Qxa5 bxa5 15. Rxb8}. If the bishop were not
        under attack, he would stand a fair chance of meeting Black's attack.

        After \variation[invar]{14. gxf4 exf4 15. Bxf4 Qxe2} Black has some extra material
        even though the attack has fizzled.
    }
    
    \mainline[level=1]{13...f4 14.gxf4 }
    
    \chessboard
    
    \mythoughts{
        When I calculated a move ago, I thought capturing on f4 was my only try. Now it feels weak:
        Black no longer has a real attack.

        \variation[invar]{
            14... Qh5
        } is both forcing and in line with my plan to sharpen the game.
    }
    
    \myreflection{
        My hesitation is obvious. One move earlier, capturing on f4 looked like my only option;
        suddenly it looks poor, even though I already sensed that the attack was gone and only
        concrete gains remained. In-between moves such as \variation{14... Qh5} are easy to
        overlook when you calculate several moves deep. They are also forcing and make the knight
        retreat. As a practical rule, leading with such moves is useful: the opponent must answer,
        and the tree stays smaller. I then see how strong \variation{10... b6} was: White's play
        stalls. Queen and rook are tied to the queenside; while they lack a breakthrough there,
        Black keeps realistic chances on the kingside. (Sometimes White gets rolling too, but Black
        may still win by mating first.)
    }

    The computer prefers \variation[invar]{14... e4}. After \variation[invar]{15. dxe4 Nxe4}
    White must react because a fork wins the exchange. Black can then snap the c4-pawn. Compared
    with my original scheme where the queen crashes through on e2, this leaves Black better placed:
    in my line White gets an open e-file and ideas against the king. Side-by-side comparisons like
    this are a handy way to choose at the board.
    
    \mainline[level=1]{14...Qh5 15.Nf3}
    
    \mythoughts{
        \variation[invar]{15... Bh3} trades a defender while drawing the queen closer to the enemy king.
    }
    \mainline[level=1]{15...Bh3
    16.Bxh3 Qxh3 17.Kh1 }
    
    \chessboard

    \mythoughts{
        The original plan was \variation[invar]{17... Ng4}. After \variation{18. Rg1} Black
        cannot continue the attack: White has \variation{19. Rxg4} with counterplay, and the g-file opens.
        \variation{18... exf4 19. Bd4} and suddenly White eyes the black king, though before Black's 17th move
        White had almost nothing active.

        I also dislike \variation[invar]{17... e4} because it hands the bishop the d4-square
        against the g7-pawn.

        I still need more weight on the kingside. The f8-rook is already part of the attack.
        White cannot grab on e5 with the f-pawn because ...Rxf3 and mate on h2 hangs in the air.
        The a8-rook is passive; ...Rbe8 intends ...Rg6 or ...Rh6 via e6.
    }

    \myreflection{
        The line of thought above is reasonable. The engine is not too harsh on \variation{17... e4} because after
        \variation{18. dxe4 Nxe4 19. Rg1 c5} Black is still alright. Part of me likes not forcing the attack
        if a won ending is on offer; on the other hand I dislike giving White the g-file.
    }
    \mainline[level=1]{17...Rbe8 18.Qb2}
    
    \chessboard

    \mythoughts{
        \variation[invar]{18... Re6} loses because of \variation{19. Rg5}.
        \variation[invar]{18... Ng4 19. Bd2 e4} Black destroys White's pawn structure.
        \variation[invar]{18... e4} needs to be calculated. \variation{19. dxe4 Rxe4} White 
        cannot play \variation{19. Rg1} because \variation{19... Rxe3} and then mate with \variation{20... Ne4}.

        I cannot tell which continuation is stronger.
    }

    \myreflection{
        After \variation{18... Ng4 19. Rg1} Black still keeps an edge. I would feel jarred:
        the attack is over. Pushing ...e4 a tempo later was a sound idea, but hoping for White's \variation{19. Bd2} was naive.
    }
      
    \myreflection{
        I get lost calculating \variation{18... e4 19. dxe4} because recapturing with the knight promises a clearly better ending:
        \variation{19. dxe4 Nxe4 20. Rg1 Nxc4 21. Qb3 Qe6 22. Rg2 Nxe3 23. fxe3 Qxb3 24. Rxb3 }. That sequence never appears on the board and Black's edge evaporates.
    }

    \myreflection{
        The practical choice needs little calculation: \variation{18... e4} is better because it forces the pace.
        After \variation{18... Ng4} I may easily overlook something and rebuild the plan from scratch.
        The whole point of \variation{18... Ng4} was to follow with \variation{19... e4}, so play the pawn break first.
    }

    \mainline[level=1]{18...e4 19.Ng5 }
    
    \mythoughts{
        White's answer is natural. After \variation[invar]{19... Qh5 20. dxe4 Ng4 21. Nf3 Rxe4}
        Black threatens ...Rf4-f3 and mate on h2. I go for it.
    }

    \myreflection{
        We will see that \variation{21... Rxe4} is not forced---Black has an even stronger option.
    }
    \mainline[level=1]{19...Qh5 20.dxe4 Ng4 21.Nf3}
    
    On move 19 I missed \variation{21... Nxc4}.
    \mainline[level=1]{21... Nxc4}
    
    \myreflection{
        Here is a pattern I know too well: I find a good plan and a good move yet miss
        a stronger one. Better moves are often easier to spot when they are forcing: the opponent must react.
        Grabbing on c4 is superior because the white queen must bail out at once. \variation[invar]{21... Rxe4 22. Qb5}
        Black probably still keeps an edge.
    }
    \mainline[level=1]{22.Qb3 Rxe4}

    Capturing on e4 was part of the plan as well.

    \mainline[level=1]{23.Rg1}
    
    \chessboard

    \mythoughts{
        \variation[invar]{
            23... Rexf4 24. Bxf4 Rxf4 25. Rxg4
        } is too committal: I see no way to press for a win once everything trades down.

        \variation[invar]{
            23... Kh8 
        } threatens ...Nxe3. \variation{24. Rxg4 Qxg4 25. Qc2} leaves h7 very tender.
    }

    \myreflection{
        I failed to find the right move. \variation{23... Kh8} is adequate, and I overlooked the game continuation.
        \variation{23... Kh8 24. Rxg4 Qxg4 25. Qc3 Nxe3 26. fxe3 Rxe3} wins for Black.
        \variation{23... Kh8 24. Rxg4 Qxg4 25. Rg1} leaves Black up material with very active pieces and realistic winning chances.
    }
    \mainline[level=1]{23...d5 24.Qb5}
    
    \chessboard

    \mythoughts{
        Black cannot improve further without landing a decisive blow---and the pieces call for one.

        \variation[invar]{
            24... Rexf4 25. Bxf4 Rxf4 26. Rxg4 Rxg4
        } I am unsure Black can convert; he is only a pawn up.

        \variation[invar]{
            24... Rxe3 25. Rxg4 Qxg4 26. Rg1 Qxf4 27. fxe3 Nxe3
        } Again Black is a pawn ahead, but White retains some activity on the g-file.

        \variation[invar]{
            24... Nxf2 25. Bxf2 Rxe2
        } Black is crashing through.

    }

    \myreflection{
        I believe I found something stronger than Carlsen's choice, yet my analysis is shaky.
        \variation{24... Rexf4 25. Bxf4 Rxf4 26. Qd7} wins a pawn, but with queens on the board White keeps some counterplay.

        \variation{
            24... Rxe3 25. Rxg4 Qxg4 26. Rg1 Qxf4 27. Qxd5+ Qf7 28. Qxf7+ Rxf7 29. fxe3 Nxe3 
        } After the queens come off, Black is simply winning.
    }
    \mainline[level=1]{24...Rfxf4? 25.Bxf4 Nxf2+ 26.Kg2 Rxe2 27.Kf1}
    
    What follows is no longer my chosen continuation---I had already found a win earlier. Along the actual moves Black cannot convert, for example:

    \variation[invar]{
        27... Qxf3 28. Rg3 Qe4 29. Qd7 Qxb1+ 30. Kxe2 Qd1+ 31. Kxf2 Qc2+ 32. Kg1 Qd1+ 33. Kf2 Qc2+ 34. Kg1 
    }
    \mainline[level=1]{27... Ne4}
    
    \mainline[level=1]{28.Rxg7+}
    
    Capturing the rook also leads to a draw, for example
    \variation[invar]{
        28... Kxg7 29. Qd7+ Qf7 30. Qg4+ Qg6 31. Qd7+ Kh8 32. Qd8+ Qg8 33. Qxg8+ Kxg8 34. Kxe2 Nxc3+ 35. Kd3 Nxb1 
    }
    \mainline[level=1]{28...Kf8 29.Kxe2 Nxc3+ 30.Kf2 Nxb5
    31.Rbg1 Nc3 32.Rxc7 Ne4+ 33.Ke1 Nc5 34.Rc8+ Kf7 35.Rc7+ Kf8 36.Rc8+ Kf7 37.Rc7+ Kf8 
    }

    \myreflection{
        I am glad my calculations were not wholly off.
        The Grand Prix Attack is worth keeping in mind: the ideas are fairly straightforward. Svidler was awful in the opening and middlegame and escaped only by a miracle.

        I take away that forcing moves deserve the first look. I will miss things as a human, and mimicking an engine by trying to calculate everything is pointless; human-style chess is more practical.

        When the position calls for direct play, the right move is often not hard to find. I still need study in these typical setups---chess offers plenty of them.

        I have been watching live games from the Candidates. The human commentators explain plans in plain language without drowning the audience in lines; I should annotate the same way.
    }
\end{multicols}